% page 590 du fac-simile (image p-589 du scan)
% bloc 157.5 x 246.3 mm ; marge gauche 8.0 mm
\pgc{-12.700mm}{0.000mm}{
\l{\cel{2}582}
\l{}
\l{\sou{tinio}. (\sou{patol.})Morbo di la pelo haroza, che la homo. - FISL}
\l{}
\l{\sou{tinklar}. (\sou{netrans.})\cel{2}(\sou{aludante klosho quan onu frapas sur}}
\l{\cel{3}nur un latero per la frapilo). Produktar sonuni qui inter-}
\l{\cel{3}sucedas lente. - EFI.}
\l{}
\l{\sou{tinlauro}. (\sou{bot.)}\cel{2}Arboreto ek la genero \textquotedbl{}viburni\textquotedbl{}. L.\cel{1}\sou{vibur}-}
\l{\cel{3}num.tinus.-}
\l{}
\l{\sou{tintar}. (\sou{trans}.,\cel{1}\sou{ye,}\cel{1}ad). I. Impregnar (filo, lano, texuro,}
\l{\cel{3}felo, e c.) ye\cel{2}substanco qua igas lu havar ta o ca koloro.}
\l{\cel{3}- II. Igar stofo ja tintita, ma di qua la koloro esas velk-}
\l{\cel{3}inta, aquirar koloro fresha. - DEFIS.}
\l{}
\l{\sou{tipo}.I. Peco de ligno, de metalo, e c., qua havas stampuro qua}
\l{\cel{3}esas destinita reproduktar stampuro stmkte sama. - II. Peco}
\l{\cel{3}de gisuro di qua la stampuro formacas la literi imprimala.}
\l{\cel{3}- III. Koncepto ideala di la traiti generala segun qua for-}
\l{\cel{3}macesis kozo. - DEFIRS.}
\l{}
\l{\sou{tipografar}. (\sou{trans.}) Imprimar per tipi. - DEFIRS.}
\l{}
\l{\sou{tipulo}.- (\sou{zool.}) Genero de insekti diptera, kun mikra korpo,}
\l{\cel{3}ek la familio \textquotedbl{}tipulidi\textquotedbl{}, kun longa gambi, tre dina e fra-}
\l{\cel{3}jila. - L.\cel{1}\sou{tipula}. -EFISL.}
\l{}
\l{\sou{tirar}. (\sou{trans., de, ek, per}). I. Venigar ula-direcione, per}
\l{\cel{3}prenar (la persono, la kozo) ye un ek lua parti quan onu}
\l{\cel{3}adduktas ad su : (a) esante sen-mova, (b) marchante ta-}
\l{\cel{3}direcione. - II. Adduktar (persono, kozo) ek la loko ube}
\l{\cel{3}lu esas inkluzita, retenata. - III. (\sou{pri kameno, furnelo,}}
\l{\cel{3}e\cel{1}\sou{c., alimentata per karbono o ligno}). Subisar en su ta}
\l{\cel{3}aero-fluo per qua la aero necesa por la kombusto atrakt-}
\l{\cel{3}esas a la herdo por remplasar laeero varmigita qua ek-}
\l{\cel{3}iras la tubo fluala, tranante kun su la fumuro. - FIS.}
\l{}
\l{\sou{tirado}. (\sou{teatro)}. To quon rolo recitas quaze \textquotedbl{}bloke\textquotedbl{}. - DEFS.}
\l{}
\l{\sou{tiraliar}. (\sou{netrans.})\cel{2}(\sou{armeo}).Pafar segun sua arbitrio, sen}
\l{\cel{3}vartar la komando de sua chefi, kontre la enemiko. - DF.}
\l{}
\l{\sou{tirano}. I. La persono qua, suprege-autoritatoza en la stato,}
\l{\cel{3}uzas sua povo ye maniero opresanta. - II. (\sou{metaf}.)Persono,}
\l{\cel{3}kozo, qua uzas opresante sua autoritato. - DEFIRS.}
\l{}
\l{\sou{tiritano}. Sorto de drapo, mi-lana e mi-lina. - DFS.}
\l{}
\l{\sou{tiroido}. (\sou{anat.}) Glando an la parto infra di la laringo, sur}
\l{\cel{3}la serio-unesma ringi di la trakeo. - EFIS.}
\l{}
\l{\sou{tirozimazo(biol}.)Oxidazo, en la fungi di la genero L.\cel{1}\sou{russu}-}
\l{\cel{3}\sou{la,}\cel{1}qua povas oxidigar tirozino.}
\l{}
\l{tirozino.(\sou{biol.}) HO. C6H4. CH2. CH(NH2) CO2 H, trovata en la}
\l{\cel{3}hepato, en la spleno, en la pankreato. -}
}
